\begin{proposition}[Open Nested Intervals Need Not Have Nonempty Intersection]
\label{prop:open-nested-intervals-need-not-have-nonempty-intersection}
There exists a sequence $(I_n)_{n\in\mathbb{N}}$ of nonempty bounded open
intervals in $\mathbb{R}$ such that
\[
  (\forall n\in\mathbb{N})(I_{n+1}\subseteq I_n),
\]
but
\[
  \bigcap_{n=1}^{\infty}I_n=\varnothing.
\]
One such sequence is
\[
  I_n=\left(0,\frac{1}{n}\right).
\]
\hyperref[prf:open-nested-intervals-need-not-have-nonempty-intersection]{Go to proof.}
\end{proposition}
\end{propositionbox}

\begin{remark*}[Standard quantified statement]
\[
\exists (a_n)\subseteq\mathbb{R}\;\exists (b_n)\subseteq\mathbb{R}\;
\left[
\begin{aligned}
&\forall n\in\mathbb{N}\;(a_n<b_n)
\land
\forall n\in\mathbb{N}\;((a_{n+1},b_{n+1})\subseteq (a_n,b_n))\\
&\land
\forall x\in\mathbb{R}\;\exists n\in\mathbb{N}\;(x\notin(a_n,b_n))
\end{aligned}
\right].
\]
\end{remark*}

\begin{remark*}[Predicate reading]
\[
\exists F\;\exists\mathcal{I}\;\exists C\;\bigl(
\mathcal{I}=\mathsf{IndexedFamily}(F,\mathbb{N},\mathbb{R})
\land C=\bigcap_{n=1}^{\infty}I_n
\land \Phi=\mathsf{OpenBoundedInterval}
\land
\operatorname{MembersSatisfy}(\mathcal{I},\Phi,\mathbb{R})
\land
\operatorname{Nested}(\mathcal{I},\subseteq)
\land
\operatorname{IsEmpty}(C)
\bigr).
\]
\end{remark*}

\begin{remark*}[Negated quantified statement]
\[
\forall (a_n),(b_n)\;\forall\mathcal{I}\;
\left[
\mathcal{I}=\mathsf{IndexedFamily}((a_n,b_n),\mathbb{N},\mathbb{R})
\land
\operatorname{MembersSatisfy}(\mathcal{I},\mathsf{OpenBoundedInterval},\mathbb{R})
\land
\operatorname{Nested}(\mathcal{I},\subseteq)
\Longrightarrow
\exists x\in\mathbb{R}\;\forall n\in\mathbb{N}\;(x\in(a_n,b_n))
\right].
\]
\end{remark*}

\begin{remark*}[Negation predicate reading]
\[
\forall (a_n),(b_n)\;\forall\mathcal{I}\;
\left[
\mathcal{I}=\mathsf{IndexedFamily}((a_n,b_n),\mathbb{N},\mathbb{R})
\land
\operatorname{MembersSatisfy}(\mathcal{I},\mathsf{OpenBoundedInterval},\mathbb{R})
\land
\operatorname{Nested}(\mathcal{I},\subseteq)
\Longrightarrow
\exists x\in\mathbb{R}\;\forall n\in\mathbb{N}\;(x\in(a_n,b_n))
\right].
\]

The negation predicate reading is the predicate-level form of the displayed negated statement.
\end{remark*}

\begin{remark*}[Failure modes]
\begin{description}
  \item[Exposition.]
  Failure is witnessed by the displayed negated or contrapositive condition.

  \item[\operatorname{MembersSatisfy} / \operatorname{Nested}.]
  \[
\forall (a_n),(b_n)\;\forall\mathcal{I}\;
\left[
\mathcal{I}=\mathsf{IndexedFamily}((a_n,b_n),\mathbb{N},\mathbb{R})
\land
\operatorname{MembersSatisfy}(\mathcal{I},\mathsf{OpenBoundedInterval},\mathbb{R})
\land
\operatorname{Nested}(\mathcal{I},\subseteq)
\Longrightarrow
\exists x\in\mathbb{R}\;\forall n\in\mathbb{N}\;(x\in(a_n,b_n))
\right].
\]
\[
\exists F\;\exists\mathcal{I}\;\exists C\;\bigl(
\mathcal{I}=\mathsf{IndexedFamily}(F,\mathbb{N},\mathbb{R})
\land C=\bigcap_{n=1}^{\infty}I_n
\land \Phi=\mathsf{OpenBoundedInterval}
\land
\operatorname{MembersSatisfy}(\mathcal{I},\Phi,\mathbb{R})
\land
\operatorname{Nested}(\mathcal{I},\subseteq)
\land
\operatorname{IsEmpty}(C)
\bigr).
\]
\end{description}
\end{remark*}

\begin{remark*}[Interpretation]
Open intervals can lose the limiting endpoint that every closed truncation
would retain.
\end{remark*}

\begin{dependencies}
\hyperref[def:open-interval]{Open Interval};
\hyperref[def:bounded-subset-of-r]{Bounded Subset of the Real Line}.
\end{dependencies}

